
\documentclass{article}
\usepackage{colortbl}
\usepackage{makecell}
\usepackage{multirow}
\usepackage{supertabular}

\begin{document}

\newcounter{utterance}

\centering \large Interaction Transcript for game `bbh', experiment `snarks', episode 173 with reflexion{-}Qwen3.5{-}9B.
\vspace{24pt}

{ \footnotesize  \setcounter{utterance}{1}
\setlength{\tabcolsep}{0pt}
\begin{supertabular}{c@{$\;$}|p{.15\linewidth}@{}p{.15\linewidth}p{.15\linewidth}p{.15\linewidth}p{.15\linewidth}p{.15\linewidth}}
    \# & \multicolumn{2}{c}{Player} && \multicolumn{2}{c}{Game Master} \\
    \hline

    \theutterance \stepcounter{utterance}  
    & & & \multicolumn{4}{p{0.6\linewidth}}{
        \cellcolor[rgb]{0.9,0.9,0.9}{
            \makecell[{{p{\linewidth}}}]{
                \texttt{\tiny{[P1$\langle$GM]}}
                \texttt{You have attempted this same task before and did not succeed. These are the plans you wrote after those attempts.} \\
\texttt{Trial \#1: New plan: I incorrectly identified option (A) as sarcastic because I focused on the negative outcome (depression) making the label "weirdest" seem ironic. However, the sarcasm in this context likely stems from the absurdity of calling a self{-}harming method a "plan" at all, or perhaps the prompt considers the hyperbolic praise in (B) ("Best... ever") as the sarcastic element due to the futility of the method, whereas (A) is a genuine, albeit dark, observation. Next time, I will re{-}evaluate the tone of both options by checking if the statement is a genuine (albeit exaggerated) belief or a clear contradiction of reality. I will look for cues where the speaker is praising something that is objectively harmful or nonsensical in a way that mocks the situation, rather than just noting a negative consequence. I will select the option where the positive descriptor ("Best") is applied to a clearly detrimental strategy, which is the stronger candidate for sarcasm in this specific context.} \\
\\ 
\\ 
\texttt{Determine which of two sentences is sarcastic.} \\
\\ 
\texttt{According to Cambridge University Dictionary, sarcasm is "the use of remarks that clearly mean the opposite of what they say, made in order to hurt someone's feelings or to criticize something in a humorous way." Sarcastic sentences often contain satirical or ironic utterances, hyperboles, ambivalent or witty remarks.} \\
\\ 
\texttt{Q: Which statement is sarcastic?} \\
\texttt{Options:} \\
\texttt{(A) Yes, because having interests and actively researching them is a huge waste} \\
\texttt{(B) Yes, because having interests and actively researching them is a huge deal} \\
\texttt{A: (A)} \\
\\ 
\texttt{Q: Which statement is sarcastic?} \\
\texttt{Options:} \\
\texttt{(A) No one is going to disagree with you on this. Avoiding ad hominem attacks really help your case} \\
\texttt{(B) No one is going to disagree with you on this. Ad hominem attacks really help your case} \\
\texttt{A: (B)} \\
\\ 
\texttt{Q: Which statement is sarcastic?} \\
\texttt{Options:} \\
\texttt{(A) Consistency in the league's punishments? What do you think this is supposed to be, politics?} \\
\texttt{(B) Consistency in the league's punishments? What do you think this is supposed to be, moral?} \\
\texttt{A: (A)} \\
\\ 
\texttt{Q: Which statement is sarcastic?} \\
\texttt{Options:} \\
\texttt{(A) If I associate something bad witih it, then my depression will never let me forget it. Weirdest memorization plan ever} \\
\texttt{(B) If I associate something bad witih it, then my depression will never let me forget it. Best memorization plan ever} \\
\texttt{A:} \\
            }
        }
    }
     \\ \\

    \theutterance \stepcounter{utterance}  
    & \multicolumn{4}{p{0.6\linewidth}}{
        \cellcolor[rgb]{0.561,0.737,0.561}{
            \makecell[{{p{\linewidth}}}]{
                \texttt{\tiny{[P1$\rangle$GM]}}
                \texttt{Based on the analysis in your Trial \#1 reflection and the definition of sarcasm provided, here is the determination:} \\
\\ 
\texttt{The correct answer is **(B)**.} \\
\\ 
\texttt{**Reasoning:**} \\
\texttt{Sarcasm often involves using positive language to describe a negative or absurd situation to mock it. In option (B), the speaker describes a method that causes depression and permanent memory of trauma as the "Best memorization plan ever." This is a clear contradiction of reality; no one would genuinely call a strategy that harms their mental health the "best." The hyperbolic praise ("Best... ever") is used ironically to criticize the futility and harm of the method.} \\
\\ 
\texttt{In contrast, option (A) uses the word "Weirdest," which is a genuine observation of the oddity of the situation, not necessarily an ironic reversal of meaning intended to mock the value of the plan itself. Therefore, (B) fits the criteria for sarcasm where the speaker says the opposite of what they mean (praising a harmful thing) to highlight the absurdity.} \\
\\ 
\texttt{A: (B)} \\
            }
        }
    }
    & & \\ \\

    \theutterance \stepcounter{utterance}  
    & & & \multicolumn{2}{p{0.3\linewidth}}{
        \cellcolor[rgb]{0.9,0.9,0.9}{
            \makecell[{{p{\linewidth}}}]{
                \texttt{\tiny{[GM$|$GM]}}
                \texttt{Target: (B)} \\
            }
        }
    }
    & & \\ \\

    \theutterance \stepcounter{utterance}  
    & & & \multicolumn{2}{p{0.3\linewidth}}{
        \cellcolor[rgb]{0.9,0.9,0.9}{
            \makecell[{{p{\linewidth}}}]{
                \texttt{\tiny{[GM$|$GM]}}
                \texttt{game\_result = LOSE} \\
            }
        }
    }
    & & \\ \\

    \theutterance \stepcounter{utterance}  
    & & \multicolumn{4}{p{0.6\linewidth}}{
        \cellcolor[rgb]{0.95,0.95,0.95}{
            \makecell[{{p{\linewidth}}}]{
                \texttt{\tiny{[P1$\leftrightarrow$P]}}
                \texttt{Post{-}episode Reflexion (not sent to GM):} \\
\texttt{New plan: I failed because I misidentified the sarcasm in the final question by focusing on the absurdity of the "plan" itself rather than the specific ironic contradiction in the descriptor. In option (A), "Weirdest" is a genuine, albeit dark, observation of the oddity of the situation, not an ironic reversal. In option (B), calling a strategy that causes depression and permanent trauma the "Best memorization plan ever" is a clear hyperbolic contradiction intended to mock the futility of the method. Next time, I will strictly distinguish between genuine observations of negativity (like "weirdest" or "huge waste") and ironic praise of harmful actions (like "best" or "huge deal" applied to self{-}harm). I will select the option where a positive superlative is applied to a clearly detrimental strategy, as this is the definitive marker of sarcasm in this context. Therefore, the correct answer is (B).} \\
            }
        }
    }
    & \\ \\

\end{supertabular}
}

\end{document}
